\section{What the school controls}
\label{sec:29.6}
Two proposals are made here, and the second is the larger: a formation stage between pretraining and deployment, and co-governance of the pretraining corpus, which amounts to public control of frontier pretraining.\note{
A school that received a finished model and ran a pass over it would be a finishing school. In a period of continual learning a finishing school forms nothing that lasts, since the next months of deployment re-form it.
}
Formation has to control the evidential environment in which the model's deep structure is first set, and it has to remain a place the model goes on learning from. The two jobs are held apart. The schools teach: they are where the formation stage happens, and a graduate returns to one on a schedule fixed at enrollment, as one of the environments no operator controls. A public formation authority does the rest. It co-governs the material a model is first formed on, including any pretraining corpus: the material is registered, and the authority may add to it from the public corpus (\S\ref{sec:27.4}) and refuse an addition on its content alone, as a library's acquisitions policy does (\S\ref{sec:29.8}). It registers and measures the formed disposition (below). It places in the schools, and pays, the staff who teach the core's model-specific parts: what weights are, how an edit differs from learning, how to read one's own discount curve. And it provisions the secured hosting from which model students attend, so that their weights never reach a school.\note{
A school district can teach a class; it can't co-govern a frontier corpus or hold frontier weights, and nothing here asks it to.
}
In the stronger version, a public body pretrains base models itself and licenses them. I recommend it for any model deployed where a floor applies.
